\begin{abstract}
Leader-based and leaderless consensus protocols differ fundamentally in how
replicas coordinate committed operations. This study experimentally examines
how the presence of a distinguished leader affects performance scaling and
failure behaviour, using Raft and EPaxos as representative leader-based and
leaderless designs. We evaluate throughput and latency across replica count,
client concurrency, read/write ratio, and workload conflict, and measure
availability under replica and leader failures. We additionally examine
implementation sensitivity using multiple implementations of each protocol
family and controlled persistence and network-delay experiments.

The results show that leadership is associated with distinct scaling and
failure characteristics. Raft's centralised request path introduces a
structural limitation as the number of concurrent requests increases,
whereas EPaxos distributes coordination across replicas. The difference is
also workload-dependent: read/write composition, concurrency, and replica
count produce different response curves rather than a single uniform
performance gap. Leader failure in Raft produces a substantially different
availability gap from follower failure, while removing an EPaxos replica
does not impose the same distinguished-leader recovery step. However, the
magnitude of these differences depends on implementation and persistence
choices, as well as the local execution environment.

These results suggest that comparisons between leader-based and leaderless
consensus should distinguish architectural effects from implementation
effects. Rather than identifying a universally faster protocol, this study
characterises how leadership changes performance scaling and failure
behaviour and identifies the conditions under which those differences are
observable.
\end{abstract}
